Hallucination detection in large language models requires uncertainty signals that capture distinct failure modes --- yet whether the two most practical signal families (consistency-based Semantic Entropy and token-level minimum log-probability) are statistically independent had not been directly tested. We measure this independence on TriviaQA dev with Llama-3.1-8B-Instruct across $N=2500$ prompts and find Pearson $|r|(\text{SE}_{N5}, \text{min\_logprob}) = 0.026$ --- near-orthogonal, far below the pre-registered collinearity threshold of 0.70. The key insight is that SE aggregates uncertainty over semantic equivalence class masses from $N=5$ stochastic samples, while min\_logprob extracts the weakest-link token confidence from a single greedy decode; these algebraically distinct operations manifest as empirically distinct signals. We further demonstrate that a cross-model LM-judge design (Qwen-2.5-7B evaluating Llama-3.1-8B outputs) produces low circularity (Spearman $\rho = -0.026$ between SE and judge labels, a numerically distinct statistic that happens to share the same rounded magnitude as the Pearson $r$, reflecting the near-zero correlation structure of this evaluation). Critically, independence is necessary but not sufficient for ensemble benefit: our conditional analysis shows SE adds no linear predictive power beyond min\_logprob (partial $R^2=0.0005$, LRT $p=0.442$ at $N=2500$). These results establish the statistical independence precondition for combining SE and min\_logprob in an ensemble, while showing that linear combinations yield no measurable gain on this benchmark --- informing that nonlinear combination methods warrant investigation in follow-up work.
